% ======================================================================
% ddG EXTENSION for the AbAffinity paper.
% Slots into the submitted "Antibody-Antigen Affinity Prediction with Chain-Aware
% Protein Language Modeling" paper: add the Method subsection after Section 2,
% the data paragraph into Section 3, and the "Section 6: Mutational ddG" block
% into Results. Compilable standalone via the fallback preamble below; to merge,
% delete the preamble/begin/end and paste the \section blocks.
% ======================================================================
\documentclass{article}
\usepackage[letterpaper,textwidth=6.0in,textheight=9in,top=1in,left=1.25in]{geometry}
\usepackage{times}
\usepackage[utf8]{inputenc}\usepackage[T1]{fontenc}
\usepackage{amsmath,amssymb,graphicx,booktabs,xcolor}
\usepackage[hidelinks]{hyperref}\usepackage{caption}
\captionsetup{font=small,labelfont=bf}
\newcommand{\ddg}{\ensuremath{\Delta\Delta G}}
\newcommand{\pkd}{\ensuremath{pK_d}}
\graphicspath{{./}}
\title{\vspace{-2em}\bfseries AbAffinity for Mutational \ddg{}: Extension Sections}
\author{}\date{}
\begin{document}\maketitle\vspace{-3em}

% ================= METHOD (append to Section 2) =================
\section{Method: repurposing AbAffinity for \ddg{}}
AbAffinity predicts absolute affinity by mapping an antibody context vector $a$ and an antigen context
vector $c$ to a cosine similarity $\hat s=\cos(a,c)$ and then to \pkd{} through its native calibration
$\pkd = {\pkd}_{\min} + ({\pkd}_{\max}-{\pkd}_{\min})\,(\hat s + 1)/2$ (Section~2). Because the change in binding
free energy upon mutation is, to first order, the difference of two absolute affinities, we repurpose the
\emph{trained, frozen} AbAffinity encoder as a physics-style \emph{affinity anchor}: we score the wild-type
and mutant complex, convert each cosine to \pkd{} with AbAffinity's own formula, and take the difference,
\begin{equation}
\widehat{\ddg} \;=\; \underbrace{s\big({\pkd}_{\mathrm{wt}}-{\pkd}_{\mathrm{mut}}\big)}_{\text{affinity anchor}}
\;+\; h\big(\phi_{\mathrm{wt}}-\phi_{\mathrm{mut}}\big),
\qquad
{\pkd}_\bullet = {\pkd}_{\min}+({\pkd}_{\max}-{\pkd}_{\min})\tfrac{\hat s_\bullet+1}{2},
\end{equation}
where $s$ is a learned scale initialized to $RT\ln 10 = 1.364\ \mathrm{kcal\,mol^{-1}}$ per \pkd{} unit (the
\pkd{}$\to$\ddg{} conversion), $\phi$ are antibody$\leftrightarrow$antigen interface-attended difference
features, and $h$ is a small gated correction network driven by mutation descriptors (log-likelihood-ratio,
physicochemical and CDR flags) that captures residual mutation-specific effects. Only the anchor scale, the
interaction, and the correction parameters are trained; ESM-2 and the AbAffinity backbone remain frozen, so
a full evaluation costs minutes on one GPU. Using AbAffinity's native \pkd{} conversion inside the anchor
(rather than differencing the raw cosine) makes the anchor a genuine predicted-\pkd{} difference and, as an
ablation (Fig.~\ref{fig:ddg}d), yields accuracy statistically identical to the cosine form ($\le 0.02$
across splits) with lower seed variance. All reported models are trained with a weighted-Huber(hinge) loss
(weight $w_k=1+\alpha\max(0,|y_k|-\tau)$, $\delta{=}1,\tau{=}2,\alpha{=}1$) plus a pairwise ranking term
(coefficient $0.5$), with reverse-mutation augmentation for antisymmetry.

% ================= EXPERIMENTAL DESIGN (append to Section 3) =================
\section{Experimental design (\ddg{})}
\textbf{Datasets and splits.} We evaluate on antibody--antigen \ddg{}: \textbf{S1131} (1{,}130 SKEMPI-derived
mutations, 112 complexes), the AB-Bind sets \textbf{AB645}/\textbf{AB1101} (473/603 recovered mutations,
19/20 complexes; AB645 is the single-point subset of AB1101), and, for out-of-domain transfer, general
protein--protein \textbf{SKEMPI 2.0} (4{,}170 recoverable entries / 280 complexes). \ddg{} is defined as
$\ddg=\Delta G_{\mathrm{mut}}-\Delta G_{\mathrm{wt}}$ in $\mathrm{kcal\,mol^{-1}}$ (positive $=$
destabilizing). For every benchmark we build a ladder of leakage-controlled 5/10-fold splits:
\emph{random} (record-level), \emph{complex-disjoint} (\texttt{GroupKFold} by PDB), and
\emph{antigen-disjoint} (\texttt{GroupKFold} by antigen sequence, the strictest, so neither antibody nor
antigen appears in both partitions). All \ddg{} numbers are mean $\pm$ s.d.\ over three seeds with per-fold
$z$-scoring; we report Pearson $r$, Spearman $\rho$, RMSE, calibrated RMSE/MAE (after a univariate linear
fit), AUROC (treating $\ddg<0$ as the affinity-enhancing class), direction accuracy, and a
per-mutation-count stratification.

\textbf{Baselines.} On \emph{identical folds} we run simple regressors on frozen ESM-2 (650M) features of
the wild-type and mutant sequences ($[\text{emb}_{\mathrm{wt}},\text{emb}_{\mathrm{mut}},
\text{emb}_{\mathrm{mut}}-\text{emb}_{\mathrm{wt}}]$): Ridge, RandomForest, XGBoost, and a small MLP; and
parameter-matched sequence modules (MLP/CNN/LSTM/lightweight Transformer over the six chain tokens). Where
protocols align we compare to published sequence predictors (AttABseq, ESM-2 zero-shot, fine-tuned
ProtAttBA) and structure predictors (DDGPred; the CIR-DDG backbones).

% ================= RESULTS =================
\section{Results: mutational \ddg{}}

\begin{figure*}[t]\centering
\includegraphics[width=\textwidth]{figure_ddg_results.pdf}
\caption{\textbf{AbAffinity-\ddg{} results} (S1131; $z$-pooled Pearson; error bars are 3-seed s.d.\ on
AbAffinity only). \textbf{(a,b)} On the complex-disjoint and antigen-disjoint splits (identical fold
columns for every method), AbAffinity-\ddg{} outperforms every frozen-ESM-2 regressor (XGBoost,
RandomForest, MLP) and the parameter-matched LSTM. \textbf{(c)} On S1131 random 10-fold CV it also leads
the published sequence methods on the same 10-fold CV (AttABseq $0.66$, DeepEP-PPI $0.41$, TransPPI $0.38$,
PIPR $0.33$, BeAtMuSiC $0.29$). \textbf{(d)} Anchor formulation: AbAffinity's native \pkd{} conversion vs
the raw cosine difference --- statistically identical. \textbf{(e)} Contribution of the learned correction:
the affinity anchor alone vs the full model. \textbf{(f)} Among large-effect mutations
($|\ddg|>2$~kcal/mol) direction accuracy and AUROC for detecting affinity-enhancing substitutions stay at
$0.94$--$0.98$.}
\label{fig:ddg}
\end{figure*}

\subsection{Chain-aware affinity generalizes to \ddg{} on unseen antigens}
On S1131 the repurposed model reaches Pearson $r=0.801\pm0.004$ (random), $0.713\pm0.017$
(complex-disjoint) and $0.675\pm0.019$ (antigen-disjoint) (Fig.~\ref{fig:ddg}a, Table~\ref{tab:ddg}). The
comparison to frozen-ESM-2 regressors reveals a clear pattern: on the \emph{random} split the embedding
regressors are as strong or stronger (XGBoost $0.851$), but \textbf{as the split becomes stricter the
AbAffinity anchor increasingly wins}, dominating on the antigen-disjoint split ($0.675$ vs XGBoost
$0.458$, $+0.22$). Parameter-matched CNN/LSTM/Transformer modules collapse under complex- and
antigen-disjoint splits ($\le 0.32$), whereas AbAffinity holds. Thus the affinity prior contributes
generalization to \emph{novel antigens} that raw embeddings and equal-capacity sequence modules do not.

\subsection{Comparison to sequence and structure predictors}
On the S1131 random benchmark (Fig.~\ref{fig:ddg}c) AbAffinity-\ddg{} ($r{=}0.801$) exceeds the frozen
sequence baselines AttABseq ($0.66$) and ESM-2 zero-shot ($0.38$); the fine-tuned ProtAttBA ($0.84$) and
structure-based DDGPred ($0.92$) are higher, as expected for a fine-tuned encoder and a structure model.
On the four external \emph{affinity} benchmarks (Section~4.1) AbAffinity already improves Pearson and RMSE
over MVSF-AB on SAbDab, AB-Bind, SKEMPI~2.0 and the held-out set. We report where fair, same-data
comparisons favour us; we do \emph{not} claim to beat supervised structure predictors on the SKEMPI
antibody cohort (see limitations).

\subsection{Out-of-domain transfer and data efficiency}
Applied to \emph{general protein--protein} SKEMPI~2.0 (out of the antibody domain), AbAffinity-\ddg{}
still attains $r=0.541\pm0.004$ complex-disjoint and $0.786\pm0.002$ random over 4{,}170 entries
(Table~\ref{tab:ddg}), i.e.\ the antibody-trained affinity representation transfers to PPI \ddg{}
and is competitive with structure-based methods' reported range. Consistent with the affinity paper's
maturation experiments, few-shot fine-tuning of the anchor (pretrained on S1131) on AB1101 rises from
$r\approx0.30$ (10\% labels) to $0.45$ (20\%), showing the same target-specific calibration behaviour for
\ddg{} ranking.

\begin{table}[t]\centering
\caption{AbAffinity-\ddg{} across benchmarks and splits (Pearson $r$, mean $\pm$ s.d.\ over 3 seeds).
XGB is the strongest frozen-ESM-2 baseline on identical folds.}
\label{tab:ddg}
\small
\begin{tabular}{lllcc}
\toprule
Benchmark & Split & AbAffinity-\ddg{} & XGB/ESM-2 & AUROC\\
\midrule
S1131 (antibody)   & random           & $0.801\pm0.004$ & $0.851\pm0.001$ & 0.880\\
S1131 (antibody)   & complex-disjoint & $\mathbf{0.713\pm0.017}$ & $0.580\pm0.005$ & 0.865\\
S1131 (antibody)   & antigen-disjoint & $\mathbf{0.675\pm0.019}$ & $0.458\pm0.004$ & 0.836\\
SKEMPI 2.0 (PPI)   & complex-disjoint & $0.541\pm0.004$ & --- & 0.678\\
SKEMPI 2.0 (PPI)   & random           & $0.786\pm0.002$ & --- & 0.775\\
AB1101 (AB-Bind)   & complex-disjoint & $0.222\pm0.027$ & 0.260 & ---\\
\bottomrule
\end{tabular}
\end{table}

\subsection{Ablation: anchor formulation and reliability}
The affinity anchor uses AbAffinity's native \pkd{} conversion (Section~4); replacing it with the raw
cosine difference (dropping the affine \pkd{} rescaling) leaves accuracy statistically unchanged
(Fig.~\ref{fig:ddg}d, identical folds): $0.801$ vs $0.803$ (random), $0.713$ vs $0.700$ (complex-disjoint),
$0.675$ vs $0.670$ (antigen-disjoint), the \pkd{} form carrying lower seed variance. This is expected: the \pkd{} conversion
is affine, so its offset cancels in the wild-type$-$mutant difference and only the scale (absorbed by $s$)
differs. The learned correction contributes beyond the anchor (Fig.~\ref{fig:ddg}e): the affinity anchor
alone reaches $0.72$/$0.69$/$0.65$ (random/complex-/antigen-disjoint), and the full model lifts these to
$0.80$/$0.71$/$0.68$. Reliability is strong on the strictest splits: among large-effect mutations
($|\ddg|>2$~kcal/mol) the model flags affinity-enhancing substitutions with AUROC $0.94$--$0.98$ and
direction accuracy $0.94$--$0.97$ (Fig.~\ref{fig:ddg}f) --- the design-relevant regime. Replacing the weighted-Huber(hinge)$+$ranking loss
with plain MSE changes $r$ by $\le 0.01$.

\begin{figure}[t]\centering
\includegraphics[width=0.62\textwidth]{figure_crossdata.pdf}
\caption{\textbf{Cross-dataset generalization.} AbAffinity-\ddg{} Pearson $r$ on three antibody \ddg{}
benchmarks and, for S1131, across all three splits (computed here for S1131; AB1101/AB645 from their
in-distribution CV). Grey cells mark splits not released for that benchmark.}
\label{fig:crossdata}
\end{figure}

\subsection{Limitations (\ddg{})}
Three limitations bound the claims. (i) On AB-Bind (AB645/AB1101; only $\sim$20 complexes) all methods
collapse under complex-/antigen-disjoint splits, so these benchmarks cannot discriminate \ddg{} methods
under leakage control. (ii) On the SKEMPI antibody interface cohort of CIR-DDG, supervised
\emph{structure} backbones (RDE-Network, DiffAffinity, Pythia-PPI; $r=0.50$--$0.61$) outperform our
sequence-only model ($r\approx0.32$ on the recoverable subset); AbAffinity beats only the zero-shot
inverse-folding baselines. (iii) On synthetic FlexddG \ddg{} (Rosetta-computed labels) the model plateaus
at $r\approx0.36$ regardless of training fraction, because the labels are a function of 3D structure a
sequence model cannot observe; this is a property of the benchmark, not of the model. In sum,
AbAffinity-\ddg{} is a structure-free predictor whose real strength is generalization to unseen antigens
and transfer to PPI, not a replacement for supervised structure \ddg{} models.

\end{document}
